\section{Conclusion}
This research introduced an efficient and scalable framework for user pairing in power-domain NOMA systems by integrating graph-based optimization with deep learning techniques. Using the standardized 3GPP TR 38.901 UMa channel model, the proposed GraphSAGE-based NOMA-GNN model effectively learned pairing strategies that closely approximated optimal graph-matching outcomes while maintaining real-time computational efficiency. The approach achieved 96.5\% of optimal Blossom throughput (636.15 Mbps vs 659.07 Mbps) with $O(N\log N)$ complexity, corresponding to 40.4$\times$ speedup compared to conventional matching algorithms. The lightweight architecture (166K parameters, $<$1 MB) is well-suited for edge base station deployment in large-scale NOMA scenarios involving hundreds to thousands of users per cell.

From a technical perspective, the developed framework bridges the gap between combinatorial optimization and data-driven inference models. By embedding physical-layer properties (SIC constraints, angular separation) directly into the graph representation, NOMA-GNN maintains interpretability while preserving wireless system consistency. The method demonstrated strong adaptability across varying user densities (N=500, 1000, 2000), maintaining 95.7-96.5\% throughput efficiency with speedup factors from 28.1$\times$ to 114.4$\times$ compared to optimal bipartite matching. Fairness analysis (Jain's index 0.9851) confirmed that the learned pairing policy ensures equitable resource allocation without sacrificing throughput. Importantly, near-optimal performance was achieved without centralized optimization or exhaustive search, demonstrating that graph neural networks can substantially enhance spectral and computational efficiency in NOMA systems.

Future research directions include: (i) extending the framework to handle imperfect CSI and channel estimation errors; (ii) developing federated learning approaches for coordinated multi-cell NOMA networks to manage inter-cell interference; (iii) investigating dynamic hybrid NOMA-OMA mode switching based on network load; (iv) incorporating mobility patterns and temporal dynamics into the graph representation; and (v) validating the approach under diverse channel models (UMi, InH) and real-world deployment scenarios. This work establishes a strong foundation for intelligent user clustering in dense 5G and emerging 6G networks, positioning Graph Neural Networks as a key enabler for autonomous, low-latency, and energy-efficient radio resource management in beyond-5G systems.
